\documentclass[10pt,a4paper]{amsart}
\usepackage[margin=19mm]{geometry}
\usepackage{amsmath,amssymb,mathtools}
\usepackage[scr=rsfs,cal=euler]{mathalfa}
\usepackage{xfrac}
\usepackage[unicode,hidelinks,bookmarks=false]{hyperref}
\newcommand{\ie}{i.e.\ }
\newcommand{\cf}{cf.\ }
\newcommand{\resp}{respectively\ }
\newcommand{\Subsection}{\subsection}
\providecommand{\nobreakdash}{\nobreak\hbox{-}}
\setlength{\emergencystretch}{2em}
\multlinegap0pt
\newtheorem{comment}{Comment}
\usepackage{amssymb}
\usepackage{bm}
\usepackage{float}
\newtheorem{definition}{Definition}
\newtheorem{lemma}{Lemma}
\newtheorem{proposition}{Proposition}
\title{Decay estimates for the two-dimensional Beam equation with potentials}
\author{Shuangshuang Chen, Han Cheng, Zijun Wan and Xiaohua Yao\textsuperscript{\dag}}
\date{Source: arXiv:2606.16793v2 (CC BY 4.0)}
\begin{document}
\maketitle

\noindent\emph{Source and licence.} Decay estimates for the two-dimensional Beam equation with potentials. Version arXiv:2606.16793v2 (CC BY 4.0), distributed by arXiv under the Creative Commons Attribution 4.0 International licence, http://creativecommons.org/licenses/by/4.0/, as recorded in the arXiv licence metadata for this exact version and shown on the abstract page https://arxiv.org/abs/2606.16793v2. Full licence notice: \texttt{licenses/arxiv-2606.16793-CC-BY-4.0.txt}.\par\smallskip

\noindent\textit{Editorial scope.} Counted unit: the article's proposition prop\_second\_2\_3 (source0.tex lines 2559--2568). The article's complete original proof of it is delivered with it (source0.tex lines 2570--2956, one \texttt{proof} environment of 387 lines). No mathematics is written, repaired or re-derived: the statement and the proof are the article's own lines, in the article's own notation, and no retained line is substituted or re-flowed.  Retained uncounted as the source-local context the proof needs: lines 1184--1193; lines 1209--1213; lines 1669--1712; lines 1717--1756; lines 2147--2153.  Nothing was omitted from any retained span, so the package carries no omission record.  No retained line contains a citation, so the wrapper carries no bibliography and no bibliography entry.  No retained line contains a figure environment, an image-inclusion command or drawing code, so no image is delivered and the package declares no asset dependency.  Editorial formatting only, in addition: an AMS \texttt{amsart} preamble that restates the article's own declarations and macros used by the retained lines, namely the environments \texttt{definition}, \texttt{lemma}, \texttt{proposition} on separate counters, together with the standard packages those lines need.  The retained environments print in this document as Definition 1, Lemma 1, Proposition 1 in order of appearance, whereas the article prints the same items with the numbering of its own full document.  The complete native source of the article as distributed in the arXiv e-print for this exact version is bundled unchanged as \texttt{sources/202966/source0.tex}.
\begin{definition}\label{defQ}
Define
$$
	Q_j := S_{j-1} - S_j \quad \text{for } j = 0,1,\dots,6, \quad \text{and} \quad Q_4^0 = S_3 - S_4^0, \quad Q_4^1 = S_4^0 - S_4.
$$
and further decompose $Q_4^1$ into the sum of the following projectors:
$$
Q_{41}^1 := S_4^0 - S_4^1, \qquad Q_{42}^1 := S_4^1 - S_4. 
$$
\end{definition}

\begin{equation}\label{condition}
	\mu> \begin{cases}11, & \mathbf{k}=0,1, \\
14, &  \mathbf{k}=2,\\
18, & \mathbf{k}=3,4. \end{cases}
		\end{equation}

\begin{lemma}\label{lemma_projection}
	Let $Q_\alpha$ $(\alpha\in\mathbb{N}_0, 0\le\alpha\le6)$ and $Q_4^0, Q_4^1, Q_{41}^1, Q_{42}^1$ be defined by \eqref{defQ}. Assume $\left|V(x)\right| \lesssim \left \langle  x\right \rangle^{-\mu}$ for some $\mu>0$ satisfying the condition \eqref{condition}. Then, for $0<\lambda\ll1$, the integral kernels admit the following decompositions:
	\begin{enumerate}
		\item[(i)] For $0\le\alpha\le6$,
		\[
		\bigl(Q_\alpha v R_0^\pm(\lambda^4)\bigr)(x,y)=\lambda^{-1-\delta_{\alpha0}}\,\Omega_{\alpha,\pm}(\lambda,x,y),
		\]
		where $\delta_{\alpha0}=1$ if $\alpha=0$ and $0$ otherwise.
		
		\item[(ii)] For $2\le\alpha\le6$, 
		\begin{align*}
		\bigl(Q_\alpha v R_0^\pm(\lambda^4)\bigr)(x,y)&=\mathcal{J}_{\alpha, \pm} (\lambda,x,y).
		\end{align*}
		\item[(iii)] For $2\le\alpha\le6$,
		setting $m_\alpha=0$ for $\alpha=2,3,4$, $m_5=1$, and $m_6=2$:
		$$
			\bigl(Q_\alpha v R_0^\pm(\lambda^4)\bigr)(x,y) =(b_0Q_\alpha v G_2^0)(x,y)+b_0(\log\lambda)\,Q_\alpha(v|\cdot|^2)(x)+ \lambda^{m_\alpha} (\mathcal{T}_{\alpha, \pm} + \mathcal{T}_\alpha)(\lambda, x, y),
        $$
		where the term $b_0 (\log\lambda) \, Q_\alpha(v|\cdot|^2)(x)$ vanishes
        for $\alpha=5,6$.
		In particular,
		\begin{align*}
\bigl(Q_4^0 v R_0^\pm(\lambda^4)\bigr)(x,y) &=(b_0Q_4^0 v G_2^0)(x,y)+b_0(\log\lambda)\,Q_4^0(v|\cdot|^2)(x)+(\mathcal{T}_{4,\pm}^0+\mathcal{T}_4^0)(\lambda,x,y),\\
\bigl(Q_4^1 v R_0^\pm(\lambda^4)\bigr)(x,y) &=(b_0Q_4^1 v G_2^0)(x,y)+(\mathcal{T}_{4,\pm}^1+\mathcal{T}_4^1)(\lambda,x,y),\\
\bigl(Q_{41}^1  v R_0^\pm(\lambda^4)\bigr)(x,y) &=(b_0Q_{41}^1  v G_2^0)(x,y)+(\mathcal{T}_{4,\pm}^{1,1}+\mathcal{T}_4^{1,1})(\lambda,x,y),\\
\bigl(Q_{42}^1  v R_0^\pm(\lambda^4)\bigr)(x,y) &=(b_0Q_{42}^1  v G_2^0)(x,y)+\lambda\bigl(\mathcal{T}_{4,\pm}^{1,2}+\mathcal{T}_4^{1,2})(\lambda,x,y).
		\end{align*}
	\end{enumerate}
	Furthermore, for $\ell=0,1,2$, these integral kernels satisfy the following bounds:
	\begin{align}\label{J}
		&\|\partial_\lambda^\ell\big(e^{\mp i\lambda|y|}\mathcal{J}_{\pm}(\lambda,\cdot,y)\big)\|_{L^2} \lesssim\lambda^{-\ell}|\log\lambda|\langle\lambda y\rangle^{-1/2},\  \text{for } \mathcal{J}_{\pm} \in \{\mathcal{J}_{\alpha,\pm}|_{2\le\alpha\le6}\},\\
&\bigl\|(b_0Q v G_2^0)(\cdot,y)\bigr\|_{L^2} \lesssim
		\begin{cases}
		\log(2+|y|), & \text{for } Q \in \{Q_\alpha|_{2 \le \alpha \le 4}, Q_4^0\}, \\[2pt]
		1, & \text{for } Q \in \{Q_4^1, Q_{41}^1, Q_{42}^1, Q_5, Q_6\},
		\end{cases} \label{lemma_projection_G} \\
&\|\partial_\lambda^\ell\mathcal{T}(\lambda,\cdot,y)\|_{L^2} \lesssim
		\begin{cases}
		\lambda^{-\ell}\log(2+|y|), & \text{for } \mathcal{T} \in \{\mathcal{T}_\alpha|_{2\le\alpha\le4}, \, \mathcal{T}_4^0\}, \\[2pt]
		\lambda^{-\ell}, & \text{for } \mathcal{T} \in \{\mathcal{T}_4^{1}, \mathcal{T}_4^{1,1}, \mathcal{T}_4^{1,2}, \mathcal{T}_5, \mathcal{T}_6\},
		\end{cases} \label{lemma_projection_T} \\   &\|\partial_\lambda^\ell\big(e^{\mp i\lambda|y|}K_{\pm}(\lambda,\cdot,y)\big)\|_{L^2} \lesssim\lambda^{-\ell}\langle\lambda y\rangle^{-1/2}, \label{lemma_projection_1}
	\end{align}
	where $K_{\pm}$ generically denotes any  term among $\Omega_{\alpha,\pm} \, (0\le\alpha\le6)$, $\mathcal{T}_{\alpha,\pm} \, (2\le\alpha\le6)$, $\mathcal{T}_{4,\pm}^j \, (j=0,1)$, or $\mathcal{T}_{4,\pm}^{1,j} \, (j=1,2)$.
\end{lemma}

\begin{lemma}\label{oscillatory}    Let \(h(x,y)\) and \(r:=r(x,y)\) be real-valued functions on \(\mathbb{R}^2 \times \mathbb{R}^2\) with \(h(x,y)>0\). Define the oscillatory integrals    
	\[
	\begin{aligned}
		\mathcal{K}^\pm(t,x,y) &= \int_0^{\infty} \cos (t \lambda^2)e^{\pm i\lambda r} \lambda \mathcal{E}^\pm(\lambda, x,y) \, d\lambda, \\
		\mathcal{N}^\pm(t,x,y) &= t^{-1}\int_0^{\infty} \sin(t\lambda^2) e^{\pm i\lambda r} \lambda^{-1} \mathcal{E}^\pm(\lambda, x,y) \, d\lambda.
	\end{aligned}
    \]  
	\begin{itemize}    
		\item[(i)] Suppose that for $\ell = 0, 1, 2$, the amplitude factor $\mathcal{E}^\pm$ satisfies    
		\begin{align}\label{condit 1}    
		\left| \partial_\lambda^{\ell} \mathcal{E}^\pm(\lambda, x,y) \right| \lesssim h(x, y) \langle \lambda r \rangle^{-\frac{1}{2}} \lambda^{-\ell}.    
		\end{align}    
		Then    
		\[    
		\left| \mathcal{K}^\pm(t,x,y) \right| + \left| \mathcal{N}^\pm(t,x,y) \right| \lesssim h(x, y) |t|^{-1}.    
		\]    
		
		\item[(ii)] Suppose that for $\ell = 0, 1, 2$, $\mathcal{E}^\pm$ exhibits a  low-frequency behavior localized by the smooth cutoff $\widetilde{\chi}_1(\lambda/2)$ supported in $[0, 2(2\lambda_0)^{1/4}]$ and    
		\begin{align}\label{condit 2}    
		\left| \partial_\lambda^{\ell} \mathcal{E}^\pm(\lambda, x,y) \right| \lesssim h(x, y) \langle \lambda r \rangle^{-\frac{1}{2}} |\log \lambda|^{-\nu} \lambda^{-\sigma - \ell} \widetilde{\chi}_1 (\lambda/2).    
		\end{align}    
		Then the following estimates hold:    
		
		$\bullet$ For $0 < \sigma < 2$ with $\nu \in \mathbb{R}$, or $\sigma = 0$ with $\nu \le 0$:    
		\[    
		\left| \mathcal{K}^\pm(t,x,y) \right| + \left| \mathcal{N}^\pm(t,x,y) \right| \lesssim \frac{h(x, y)}{\langle t \rangle^{1-\frac{\sigma}{2}} (\log(2+|t|))^{\nu}}.    
		\]    
Moreover, if $r=0,$  this estimate  remains
   valid for the larger parameter range
   $
   -2<\sigma<2$ and 
   $\nu\in\mathbb R.
   $		
		
        $\bullet$ For $\sigma = 2$ and $\nu > 1$:    
		\[    
		\left| \mathcal{K}^\pm(t,x,y) \right| + \left| \mathcal{N}^\pm(t,x,y) \right| \lesssim  \frac{h(x, y)}{\left(\log(2+|t|)\right)^{\nu-1}}.    
		\]    
	\end{itemize} 
\end{lemma}    

\begin{align}
	\label{K_1}
	& \mathcal{K}_{B}^\pm(t,x,y)=\frac{2}{\pi i}\int_0^{\infty} \cos(t \lambda^2)\lambda \widetilde{\chi}_1(\lambda) \lambda^{4-k_\alpha-k_\beta}\bigl\langle  B^\pm(\lambda)Q_\beta vR_0^\pm(\lambda^{4})(\cdot,y), Q_\alpha vR_0^\mp(\lambda^{4})(\cdot,x)\bigr\rangle \, d \lambda, \\
	\label{Lambda_1}
	& \mathcal{N}_{B}^\pm(t,x,y)=\frac{2}{\pi it}\int_0^{\infty} \sin(t\lambda^2)\lambda^{-1}\widetilde{\chi}_1(\lambda)\lambda^{4-k_\alpha-k_\beta}
	\bigl\langle  B^\pm(\lambda)Q_\beta vR_0^\pm(\lambda^{4})(\cdot,y), Q_\alpha vR_0^\mp(\lambda^{4})(\cdot,x)\bigr\rangle \, d \lambda,
\end{align}

\begin{proposition}\label{prop_second_2_3}
	Let $B^\pm(\lambda) \in \{\mathcal{M}_{4,\alpha}^{\pm}(\lambda), \mathcal{M}_{\alpha,4}^{\pm}(\lambda)\}$ with $\alpha \in \{0,1,2,3,4\}$. Then
	\[  
\sup _{x, y \in \mathbb{R}^2}\left| \mathcal{K}_{B}^\pm(t,x,y)\right| + \sup _{x, y \in \mathbb{R}^2}\left|\mathcal{N}_{B}^\pm (t,x,y)\right| \lesssim
	\begin{cases}
	\langle t\rangle^{-1}, & \alpha\in\{0,1,2,3\},\\
	\langle t\rangle^{-1}\bigl(\log(2+|t|)\bigr)^2, & {\alpha} = 4.
\end{cases}  
	\]  
\end{proposition}

\begin{proof}
	By Lemma \ref{lemma_projection}, the projected free resolvents relevant
	to the present proposition satisfy
	\begin{align}
		\bigl(Q_\alpha vR_0^\pm(\lambda^4)\bigr)(\cdot,z)
		&=
		\lambda^{-1-\delta_{\alpha0}}
		\Omega_{\alpha,\pm}(\lambda,\cdot,z),
		\label{Qalpha}\\
		\bigl(Q_4vR_0^\pm(\lambda^4)\bigr)(\cdot,z)
		&=
		\mathcal{J}_{4,\pm}(\lambda,\cdot,z),
		\label{De all Q}\\
		\bigl(Q_4^1vR_0^\pm(\lambda^4)\bigr)(\cdot,z)
		&=
		(b_0Q_4^1vG_2^0)(\cdot,z)
		+
		\mathcal{T}_4^1(\lambda,\cdot,z)
		+
		\mathcal{T}_{4,\pm}^1(\lambda,\cdot,z).
		\label{eq:resolvent_decomp-Q4}
	\end{align}
	Moreover, for $\ell=0,1,2$,
	\begin{equation}
		\begin{split}
			&	\big\|
			\partial_\lambda^\ell
			\bigl(
			e^{\mp i\lambda|z|}
			\mathcal{J}_{4,\pm}(\lambda,\cdot,z)
			\bigr) 
			\big\|_{L^2} \lesssim 
			\lambda^{-\ell} 	|\log\lambda|
			\langle\lambda z\rangle^{-1/2},\\
		&
		\big\|
		\partial_\lambda^\ell
		\bigl(
		e^{\mp i\lambda|z|}
		\Omega_{\alpha,\pm}(\lambda,\cdot,z)
		\bigr)
		\big\|_{L^2}
		+
		\big\|
		\partial_\lambda^\ell
		\bigl(
		e^{\mp i\lambda|z|}
		\mathcal{T}_{4,\pm}^1(\lambda,\cdot,z)
		\bigr)
		\big\|_{L^2}
		\lesssim
		\lambda^{-\ell}
		\langle\lambda z\rangle^{-1/2},\\
	&	\|
		\partial_\lambda^\ell
		\mathcal{T}_4^1(\lambda,\cdot,z)
		\|_{L^2}
		\lesssim
		\lambda^{-\ell},\quad 	\|(b_0Q_4^1vG_2^0)(\cdot,z)\|_{L^2}
		\lesssim1.
		\label{eq:E_bound-1}
		\end{split}
	\end{equation}
	We divide the proof of this proposition into three cases: $\alpha=0; 1\leq\alpha\leq3;$ and $\alpha=4.$
	\begin{comment}
	We shall also repeatedly use
	\begin{equation}\label{eq:weight-product}
		\langle\lambda x\rangle^{-1/2}
		\langle\lambda y\rangle^{-1/2}
		\lesssim
		\langle\lambda(|x|+|y|)\rangle^{-1/2}.
	\end{equation}
	Derivatives falling on $\widetilde{\chi}_1(\lambda)$ will always be
	absorbed into the slightly enlarged cutoff
	$\widetilde{\chi}_1(\lambda/2)$.
	\end{comment}
	\medskip
	
	
	\textbf{Case 1: $B^\pm=\mathcal{M}_{0,4}^\pm$ or
		$\mathcal{M}_{4,0}^\pm$.}
	We consider $\mathcal{M}_{0,4}^\pm$  (the symmetric case $\mathcal{M}_{4,0}^{\pm}$ follows identically). Write
	$$
	\mathcal{M}_{0,4}^\pm(\lambda)
	=
	B_0^\pm(\lambda)+B_1^\pm(\lambda),
	$$
	where
	$
	B_0^\pm(\lambda)
	=
	(\log\lambda)^{-1}
	Q_0\Lambda^\pm(\lambda)Q_4^0$
	and $
	B_1^\pm(\lambda)
	=
	Q_0\Lambda^\pm(\lambda)Q_4^1.
	$ 	By linearity, it suffices to estimate the contributions of
	$B_0^\pm$ and $B_1^\pm$ separately.
	
	We first treat $B_0^\pm$. Since
	$
	4-k_0-k_4=2,
	$
	substituting \eqref{Qalpha} and \eqref{De all Q} into both
	\eqref{K_1} and \eqref{Lambda_1} gives
	\begin{align}
		\mathcal{K}_{B_0}^\pm(t,x,y)
		={}&
		\frac{2}{\pi i}
		\int_0^\infty
		\cos(t\lambda^2)
		e^{\pm i\lambda(|x|+|y|)}
		\lambda
		\mathcal{E}_{0}^\pm(\lambda,x,y)
		\,d\lambda,
		\label{eq:K-B0-standard}\\
			\mathcal{N}_{B_0}^\pm(t,x,y)
		={}&
		\frac{2}{\pi it}
		\int_0^\infty
		\sin(t\lambda^2)
		e^{\pm i\lambda(|x|+|y|)}
		\lambda^{-1}
		\mathcal{E}_{0}^\pm(\lambda,x,y)
		\,d\lambda,
		\label{eq:N-B0-standard}
	\end{align}
	where
	\begin{align}
		\mathcal{E}_{0}^\pm(\lambda,x,y)
		={}&
		\widetilde{\chi}_1(\lambda)
		(\log\lambda)^{-1}
		\Bigl\langle
		Q_0\Lambda^\pm(\lambda)Q_4^0
		\Bigl(
		e^{\mp i\lambda|y|}
		\mathcal{J}_{4,\pm}(\lambda,\cdot,y)
		\Bigr),
		e^{\pm i\lambda|x|}
		\Omega_{0,\mp}(\lambda,\cdot,x)
		\Bigr\rangle.
		\label{eq:E-B0}
	\end{align}
	Indeed, the factor $\lambda^2$ arising from
	$\lambda^{4-k_0-k_4}$ exactly cancels the factor $\lambda^{-2}$
	in the $Q_0$-resolvent expansion.
	
	Observe that  the cosine and sine kernels have exactly the forms required in
	Lemma \ref{oscillatory}, with
	$
	r(x,y)=|x|+|y|
	$
	and the common amplitude factor $\mathcal{E}_{0}^\pm$ given by
	\eqref{eq:E-B0}.
Since $\|\partial_\lambda^\ell \Lambda^\pm(\lambda)\|_{\mathbb{B}(L^2)} \lesssim \lambda^{-\ell}$ for $\ell=0,1,2$ and the factor $(\log\lambda)^{-1}$ exactly compensates for the logarithmic growth of $\mathcal{J}_{4,\pm}$ in \eqref{eq:E-B0}, it follows from Leibniz's rule, H\"older's inequality, and \eqref{eq:E_bound-1} that
	\begin{equation*}
		%\label{eq:E-B0-bound}
		\left|
		\partial_\lambda^\ell
		\mathcal{E}_{0}^\pm(\lambda,x,y)
		\right|
		\lesssim
		\widetilde{\chi}_1(\lambda/2)
		\langle\lambda(|x|+|y|)\rangle^{-1/2}
		\lambda^{-\ell},
		\qquad
		\ell=0,1,2.
	\end{equation*}
 Lemma
	\ref{oscillatory}(ii) with $(\sigma,\nu)=(0,0)$ therefore yields
	the desired $\langle t\rangle^{-1}$ bound for both
	$\mathcal{K}_{B_0}^\pm$ and $\mathcal{N}_{B_0}^\pm$.
	
	We next treat $B_1^\pm$. By
	\eqref{eq:resolvent_decomp-Q4}, the corresponding kernels 	$\mathcal{K}_{B_1}^\pm$ and $\mathcal{N}_{B_1}^\pm$ split into
	three contributions, associated respectively with
	$Q_4^1vG_2^0$, $\mathcal{T}_4^1$, and
	$\mathcal{T}_{4,\pm}^1$. The same calculation leading to
	\eqref{eq:K-B0-standard} and \eqref{eq:N-B0-standard} shows that each
	contribution has the form required in Lemma \ref{oscillatory}. We
	therefore only specify the corresponding phases and amplitude factors.
	
	The contribution of $Q_4^1vG_2^0$ has the phase
	$
	r(x,y)=|x|
	$
	and the amplitude factor
	\begin{equation*}
		%\label{eq:E-B1-G}
		\mathcal{E}_{G}^\pm(\lambda,x,y)
		=
		\widetilde{\chi}_1(\lambda)
		\Bigl\langle
		Q_0\Lambda^\pm(\lambda)Q_4^1
	(Q_4^1vG_2^0)(\cdot,y),
		e^{\pm i\lambda|x|}
		\Omega_{0,\mp}(\lambda,\cdot,x)
		\Bigr\rangle.
	\end{equation*}
	The contribution of $\mathcal{T}_4^1$ has the same phase
	$r(x,y)=|x|$ and  the amplitude factor
	\begin{equation*}
		%\label{eq:E-B1-T}
		\mathcal{E}_{\mathcal{T}}^\pm(\lambda,x,y)
		=
		\widetilde{\chi}_1(\lambda)
		\Bigl\langle
		Q_0\Lambda^\pm(\lambda)Q_4^1
		\mathcal{T}_4^1(\lambda,\cdot,y),
		e^{\pm i\lambda|x|}
		\Omega_{0,\mp}(\lambda,\cdot,x)
		\Bigr\rangle.
	\end{equation*}
	Finally, the contribution of $\mathcal{T}_{4,\pm}^1$ has the  phase
	$
	r(x,y)=|x|+|y|
	$
	and the  amplitude factor
	\begin{align*}
		\mathcal{E}_{osc}^\pm(\lambda,x,y)
		={}&
		\widetilde{\chi}_1(\lambda)
		\Bigl\langle
		Q_0\Lambda^\pm(\lambda)Q_4^1
		\Bigl(
		e^{\mp i\lambda|y|}
		\mathcal{T}_{4,\pm}^1(\lambda,\cdot,y)
		\Bigr),
		e^{\pm i\lambda|x|}
		\Omega_{0,\mp}(\lambda,\cdot,x)
		\Bigr\rangle.
	%	\label{eq:E-B1-osc}
	\end{align*}
	
	By \eqref{eq:E_bound-1}, the derivative
	bounds for $\Lambda^\pm$, Leibniz's rule, and H\"older's inequality, each of these amplitude factors satisfies
	\begin{equation*}
		%\label{eq:E-B1-bound}
		\left|
		\partial_\lambda^\ell
		\mathcal{E}^\pm(\lambda,x,y)
		\right|
		\lesssim
		\widetilde{\chi}_1(\lambda/2)
		\langle\lambda r\rangle^{-1/2}
		\lambda^{-\ell},
		\qquad
		\ell=0,1,2,
	\end{equation*}
	where $\mathcal{E}^\pm$ denotes any of
	$\{
	\mathcal{E}_{G}^\pm,
	\mathcal{E}_{\mathcal{T}}^\pm,
	\mathcal{E}_{osc}^\pm\}
	$
	and $r$ is the corresponding phase specified above. Then applying Lemma
	\ref{oscillatory}(ii) with $(\sigma,\nu)=(0,0)$ yields  the desired $\langle t\rangle^{-1}$ bound for both
	$\mathcal{K}_{B_1}^\pm$ and $\mathcal{N}_{B_1}^\pm$.
	
	Combining the estimates for $B_0^\pm$ and $B_1^\pm$ gives
	$$
	\sup_{x,y\in \mathbb{R}^2}
	\bigl|
	\mathcal{K}_{\mathcal{M}_{0,4}}^\pm(t,x,y)
	\bigr|
	+
	\sup_{x,y\in \mathbb{R}^2}
	\bigl|
	\mathcal{N}_{\mathcal{M}_{0,4}}^\pm(t,x,y)
	\bigr|
	\lesssim
	\langle t\rangle^{-1}.
	$$

	

	\textbf{Case 2: $B^\pm=\mathcal{M}_{\alpha,4}^\pm$ or
		$\mathcal{M}_{4,\alpha}^\pm$, $1\leq\alpha\leq3$.}
	We consider $\mathcal{M}_{\alpha,4}^\pm$ (the symmetric case $\mathcal{M}_{4,\alpha}^{\pm}$ follows identically). Since
	$
	4-k_\alpha-k_4=1
	$
	and
	$
	\bigl(
	Q_\alpha vR_0^\mp(\lambda^4)
	\bigr)(\cdot,x)
	=
	\lambda^{-1}
	\Omega_{\alpha,\mp}(\lambda,\cdot,x),
	$
	the explicit factor $\lambda$ arising from
	$\lambda^{4-k_\alpha-k_4}$  in
	\eqref{K_1} and \eqref{Lambda_1} exactly cancels this
	$\lambda^{-1}$ singularity. By repeating the calculation in
	\eqref{eq:K-B0-standard}--\eqref{eq:N-B0-standard}, the two kernels $	\mathcal{K}_{\mathcal{M}_{\alpha,4}}^\pm$ and $	\mathcal{N}_{\mathcal{M}_{\alpha,4}}^\pm$ 
	take the forms required in Lemma \ref{oscillatory} with the  phase
	$
	r(x,y)=|x|+|y|
	$
	and the  amplitude factor
	\begin{align*}
		\mathcal{E}_{\alpha4}^\pm(\lambda,x,y)
		={}&
		\widetilde{\chi}_1(\lambda)
		\Bigl\langle
		\mathcal{M}_{\alpha,4}^\pm(\lambda)
		\Bigl(
		e^{\mp i\lambda|y|}
		\mathcal{J}_{4,\pm}(\lambda,\cdot,y)
		\Bigr),
		e^{\pm i\lambda|x|}
		\Omega_{\alpha,\mp}(\lambda,\cdot,x)
		\Bigr\rangle.
	%	\label{eq:E-alpha4}
	\end{align*}
	Recall that
$
\|\partial_\lambda^\ell \mathcal{M}_{\alpha, 4}^{\pm}(\lambda)\|_{\mathbb{B}(L^2)} \lesssim \lambda^{1-\ell}|\log\lambda|^{4}$ for   $\ell=0,1,2.$ 
Applying Leibniz's rule and H\"older's inequality alongside \eqref{eq:E_bound-1}, we see that the additional $\lambda$ factor from $ \mathcal{M}_{\alpha, 4}^{\pm}(\lambda)$ absorbs the logarithmic singularities:
\[  
\bigl| \partial_\lambda^{\ell} \mathcal{E}_{\alpha4}^\pm(\lambda,x,y) \bigr| \lesssim \widetilde{\chi}_1(\lambda/2) \langle \lambda r \rangle^{-1/2} \lambda^{1-\ell}|\log\lambda|^5 \lesssim \widetilde{\chi}_1(\lambda/2) \langle \lambda r \rangle^{-1/2} \lambda^{-\ell}, \quad \ell = 0,1,2. 
\]   
	Hence Lemma \ref{oscillatory}(ii) again gives
	$$
	\sup_{x,y}
	\bigl|
	\mathcal{K}_{\mathcal{M}_{\alpha,4}}^\pm(t,x,y)
	\bigr|
	+
	\sup_{x,y}
	\bigl|
	\mathcal{N}_{\mathcal{M}_{\alpha,4}}^\pm(t,x,y)
	\bigr|
	\lesssim
	\langle t\rangle^{-1}.
	$$
	
	
	\textbf{Case 3: $B^\pm=\mathcal{M}_{4,4}^\pm$.}
	This case  relies on the representations \eqref{De all Q} and \eqref{eq:resolvent_decomp-Q4}.
	Note that the expansion for $Q_4v R_0^\pm(\lambda^4)$ contains an $|\log\lambda|$ singularity, whereas the expansion for $Q_{4}^1v R_0^\pm(\lambda^4)$ does not. Moreover, none of the bounds in \eqref{eq:E_bound-1} carry the spatial weight $\omega(z)$.
	Note that  $4-k_4-k_4=0$, the intrinsic $\lambda$-factor in the definitions \eqref{K_1} and \eqref{Lambda_1} is $\lambda^0 = 1$.
Substituting  \eqref{De all Q}, \eqref{eq:resolvent_decomp-Q4} and  the  expansion 
		\begin{align}\label{eq:M44_rearranged-1}
		\mathcal{M}_{4,4}^{\pm}(\lambda) = Q_{4} \Lambda^\pm(\lambda) Q_{4} + (\log\lambda)\bigl(Q_{4} \Lambda^\pm(\lambda) Q_{4}^1+Q_{4}^1 \Lambda^\pm(\lambda) Q_{4}\bigr) + (\log\lambda)^2 Q_{4}^1 \Lambda^\pm(\lambda)Q_{4}^1,
	\end{align}
	 into \eqref{K_1} and \eqref{Lambda_1}, the kernels $\mathcal{K}_B^\pm$ and $\mathcal{N}_B^\pm$ reduce to finite linear combinations of canonical oscillatory integrals of the form:
	\begin{equation}\label{eq:canonical_oscillatory_prop1}
		\frac{2}{\pi i}	\int_0^\infty \lambda \cos(t\lambda^2) e^{\pm i\lambda r(x,y)} \mathcal{E}_{f,g}^\pm(\lambda, x, y) \, d\lambda \quad \text{and} \quad
		\frac{2}{\pi it} \int_0^\infty \lambda^{-1} \sin(t\lambda^2) e^{\pm i\lambda r(x,y)} \mathcal{E}_{f,g}^\pm(\lambda, x, y) \, d\lambda,
	\end{equation}
	where the amplitude factor $\mathcal{E}_{f,g}^\pm$ is defined via the inner product:
	\begin{equation}\label{eq:amplitude_def_prop1}
		\mathcal{E}_{f,g}^\pm(\lambda, x, y) = e^{\mp i\lambda r(x,y)} \widetilde{\chi}_1(\lambda) \big\langle M_{sub}^\pm(\lambda) f(\lambda, \cdot, y), \, g(\lambda, \cdot, x) \big\rangle.
	\end{equation}
	Here, $M_{sub}^\pm(\lambda)$ is a constituent operator from \eqref{eq:M44_rearranged-1}, while $f$ and $g$ are corresponding states selected from the expansions in \eqref{De all Q} and \eqref{eq:resolvent_decomp-Q4}. The phase function $r(x,y) \in \{0, |x|, |y|, |x|+|y|\}$ neutralizes the exponential phases embedded within $f$ and $g$.

	To uniformly bound \eqref{eq:canonical_oscillatory_prop1}, we apply H\"older's inequality to \eqref{eq:amplitude_def_prop1}, tracing the $\lambda$-logarithmic singularities inherited from each component:
	\begin{itemize}
		\item For $M_{sub}^\pm(\lambda) = Q_{4} \Lambda^\pm(\lambda) Q_{4}$, it acts between the states $f = \mathcal{J}_{4,\pm}$ and $g = \mathcal{J}_{4,\mp}$. The product of their individual $L^2$-bounds yields an amplitude factor constrained by $ |\log\lambda|^2$.
		\vskip0.2cm
		\item For $M_{sub}^\pm(\lambda) = (\log\lambda) Q_{4} \Lambda^\pm(\lambda) Q_{4}^1$ (the symmetric case $(\log\lambda) Q_{4}^1 \Lambda^\pm(\lambda)Q_{4}$ follows identically),
		the operator pairs $f\in\{b_0 Q_{4}^1 v G_2^0, \mathcal{T}_{4,\pm}^{1}, \mathcal{T}_{4}^{1}\}$ and $g = \mathcal{J}_{4,\mp}(\lambda, \cdot, x)$. The inherent logarithm in the operator multiplies the singularity of $g$, yielding a total $\lambda$-singularity of $ |\log\lambda|^2$.
		\vskip0.2cm
		\item For $M_{sub}^\pm(\lambda) = (\log\lambda)^2 Q_{4}^1 \Lambda^\pm(\lambda) Q_{4}^1$, this operator pairs with the bounded, $\log$-free $Q_{4}^1$-expansions on both sides, i.e., $f,g\in\{b_0 Q_{4}^1 v G_2^0, \mathcal{T}_{4,\pm}^{1}, \mathcal{T}_{4}^{1}\}.$  The $\lambda$-singularity arises entirely from the operator itself and is bounded by $|\log\lambda|^2$.
	\end{itemize}
	Thus, for all cross-terms, applying Leibniz's rule and H\"older's inequality to the amplitude factor yields
	\[
	\bigl| \partial_\lambda^{\ell} \mathcal{E}_{f,g}^\pm(\lambda,x,y) \bigr| \lesssim \widetilde{\chi}_1(\lambda/2) \langle \lambda r(x,y) \rangle^{-1/2} \lambda^{-\ell} |\log\lambda|^2, \quad \ell = 0,1,2.
	\]
	An application of Lemma \ref{oscillatory}(ii) with parameters $(\sigma,\nu)=(0,-2)$ then yields $$
	\sup_{x,y\in \mathbb{R}^2}
	\bigl|
	\mathcal{K}_{\mathcal{M}_{4,4}}^\pm(t,x,y)
	\bigr|
	+
	\sup_{x,y\in \mathbb{R}^2}
	\bigl|
	\mathcal{N}_{\mathcal{M}_{4,4}}^\pm(t,x,y)
	\bigr|
	\lesssim
	\langle t\rangle^{-1}\bigl(\log(2+|t|)\bigr)^2.
	$$
\end{proof}

\end{document}
